% ======================================================================
% SECTION 5: LIMITATIONS AND FUTURE WORK
% Five thematic subsections: (i) approximations vs generated vs
% executed accounting, (ii) 60min vs 1min formal deltas, (iii) cross-
% simulation limitations, (iv) Track A vs Track B future directions,
% (v) future work to reduce approximations, reduce surgery time, and
% increase time resolution.
% Bracketed prompts for future Claude Code Opus 4.7 1M Max processing.
% Do not process the bracketed prompts here.
% ======================================================================

\subsection{Approximations vs Generated vs Executed Accounting}
\label{sec:limits-accounting}

[LIMITATIONS SUBSECTION 1 PROMPT - PROCESS DURING THE PAPER POPULATION PASS]

[Document what was approximated from the upstream instructions vs
what was generated by Claude Code Opus 4.7 1M Max vs what was
executed. Render a three-way LaTeX table using
{>{\raggedright\arraybackslash}p{Xcm}} for each column with columns:
artifact category, approximated by upstream
(physical-ai-oncology-trials/competitions/instructions/one_minute_variant/),
generated by Claude Code Opus 4.7 1M Max
(kevinkawchak/robotic-surgeries/2030-gbm-1min/), executed by Claude
Code (kevinkawchak/robotic-surgeries/2030-gbm-1min/outputs/). For
every row, name the exact files. The artifact categories are: (a)
number of instruction files (upstream 12 hand-authored Markdown
files), (b) Cargo / pyproject / docker-compose toolchain (generated
in 2030-gbm-1min/; executed via build artifact in outputs/logs/),
(c) 200-channel sensor schema (upstream commit_02_sensors_1min.md;
generated in schemas/sensor_record_4arm.{schema.json, proto, avsc};
executed in outputs/sensors/), (d) 7-DOF DH parameter table
(upstream robot_specification_neurospeed.md; generated in
config/kinematics_4arm.yaml; executed in outputs/xyz_mapping/), (e)
16-iteration sweep design (upstream commit_04_iterations_1min.md;
generated in config/iterations.yaml and src/simulation/iterate_1min.py;
executed in outputs/iterations/), (f) composite score formula
(upstream commit_05_competition_1min.md; generated in
src/metrics/compute_1min.py; executed in outputs/metrics/), (g) LLM
tournament (upstream commit_05_competition_1min.md; generated in
src/llm/compare_agent_1min.py and prompts/comparison_prompt_1min.md;
executed in outputs/comparison/ and outputs/comparison_robot_vs_human/),
(h) Zenodo L0 deposition (upstream zenodo_archive_protocol.md;
generated in src/zenodo/patch_pointers.py and the pointer JSON files;
executed in outputs/iterations/*_L0_raw.zenodo_pointer.json), (i)
sensor sample table (upstream sensor_specification_10khz.md
approximates the schema; generated in src/sensors/ingest_4arm.py;
executed as the 54 columns by 1001 rows file at
outputs/sensors/sensor_sample_4arm.csv), (j) ASCII diagrams (upstream
multi_arm_coordination.md provides the heartbeat sketch; generated
through docs/architecture_overview_4arm.txt; executed as 6 curated
ASCII diagrams at outputs/diagrams/). Highlight every place where
the upstream instruction approximates a count or a size and the
downstream pass committed to a specific number; note the
approximations for future tightening.]

\subsection{60-Minute vs 1-Minute Formal Delta Accounting}
\label{sec:limits-deltas}

[LIMITATIONS SUBSECTION 2 PROMPT - PROCESS DURING THE PAPER POPULATION PASS]

[Document the formal delta between the prior 60-minute study (read
from
kevinkawchak/physical-ai-oncology-trials/competitions/instructions/
and its 60-minute subdirectory if present, plus the prior paper at
kevinkawchak/physical-ai-oncology-trials/new-trial/national-24-7-trial/paper/full-paper/final-paper/)
and the current 1-minute study (read from
kevinkawchak/physical-ai-oncology-trials/competitions/instructions/one_minute_variant/
and the entire kevinkawchak/robotic-surgeries/2030-gbm-1min/ tree).
Render a side by side LaTeX table using
{>{\raggedright\arraybackslash}p{Xcm}} for each column with columns:
dimension, prior 60 minute scope, current 1 minute scope, ratio.
Cover at minimum the following dimensions: number of arms (prior
single arm; current 4 cooperating arms), tool assignment per arm
(prior shared tools; current dedicated arm 1 hybrid u-w-p, arm 2
bipolar plus irrigation, arm 3 suction plus collection, arm 4 iMRI
plus 5-ALA plus ultrasound), sensor channels (prior single arm
channel stack at lower rate; current 50 channels per arm x 4 arms
at mixed 1 kHz commands plus 10 kHz force = 200 channels), cartesian
coordinate update rate (prior lower rate; current 1 kHz commands
plus 10 kHz force), iteration count (prior fewer iterations; current
16 deterministic), tournament entity count (prior 4 entities at the
sponsor level; current 4 entities at the robot level), structural
caveat (prior pharmacology dimension only; current time dimension
1 min vs 1 h structural), procedure duration (60 minutes vs 60
seconds; 60x compression), and committed footprint (prior
multi-thousand artifact set across sponsor and site simulations;
current approximately 8.2 MB plus 1.5 MB overhead inside the 10 MB
commit cap). Cite kawchak_2026_19994945 for the prior 60 minute
scope and repo-robotic-surgeries for the current 1 minute scope.
This subsection is the formal home of the 60min vs 1min delta
accounting; downstream prose elsewhere should cross reference here.]

\subsection{Cross-Simulation Limitations}
\label{sec:limits-cross}

[LIMITATIONS SUBSECTION 3 PROMPT - PROCESS DURING THE PAPER POPULATION PASS]

[Document the limitations that apply across all artifacts in the
2030-gbm-1min/ tree. State explicitly that (a) the synthetic patient
PAT-GBM-0001 (4.2 cm right frontal IDH-wildtype glioblastoma) is not
a real patient; no IRB, no informed consent, and no real EHR signal
was used; (b) Claude Code Opus 4.7 1M Max output is not deterministic
across runs; a re-generation of the 2030-gbm-1min/ tree may produce
different file names, sizes, and prose, although the per-iteration
Parquet aggregates under seed 20260510 are bit deterministic for
fixed parameter tuple; (c) the synthetic 16-iteration sweep has not
been validated against any real-patient cohort or TRIPOD+AI / CREMLS
reporting floor (cite Collins2024TRIPODAI, ElEmam2024CREMLS); (d)
the hypothetical 2030 Medtronic NeuroSpeed 1.0 hardware does not
exist; the comparison against Medtronic ROSA ONE Brain v3.0 (cite
rosa-one-brain, Lefranc2014ROSAStereotaxy) is paper-only; (e) the
C++20 control loop, C++20 1 kHz heartbeat layer, and Rust 2021
high-throughput runner are documented but may not have been compiled
in this generation pass (read
kevinkawchak/robotic-surgeries/2030-gbm-1min/outputs/reports/limitations.md
in full and reproduce the partial-verification status); (f) the LLM
tournament rationale at outputs/comparison/comparison.json and
outputs/comparison_robot_vs_human/comparison.json carries the
structural-time-dimension caveat in every round (1 minute robot
trivially beats 1 hour human-baseline on time), and the caveat must
not be removed downstream; (g) the work is not endorsed or sponsored
by any trial sponsor, FDA, CRO, site, IRB, regulator, or medical
society and was generated using Artificial Intelligence.]

\subsection{Track A vs Track B Future Directions}
\label{sec:limits-tracks}

[LIMITATIONS SUBSECTION 4 PROMPT - PROCESS DURING THE PAPER POPULATION PASS]

[Document the two future-work tracks that extend the computational
advantages demonstrated here into real-world applications. Track A:
single big model performing all tasks. A future generation of Claude
Code (or an equivalent foundation-model agent) holds the entire
2030-gbm-1min/ repository in context and emits all per-iteration
patient predictions, per-iteration safety signals, and per-iteration
FDA-bound signals from a single 1M token inference pass. Track A is
architecturally simple but has higher inference cost per decision
and a single point of regulatory accountability. Track B: single big
model that creates smaller agents performing smaller tasks locally.
A future generation uses the 1M token context model only for
orchestration and routing, while delegating per-task prediction
(per-arm force prediction, per-iteration composite score, per-round
tournament verdict) to smaller specialized agents that run on
commodity hardware (Core i5-6200U or similar). Track B distributes
regulatory accountability but lowers per-site inference cost and
improves the security posture for PHI handling. Render the trade
offs as a LaTeX table using {>{\raggedright\arraybackslash}p{Xcm}}
for each column with columns: property, Track A: big single model,
Track B: big model plus small agents. Include at least the following
rows: architecture, inference cost, regulatory accountability, PHI
handling, hardware floor, update cycle, extension target. The
extension target for Track A is the cloud-orchestrated 60-second
resection that points at a real NeuroSpeed-class robot; for Track B
it is the on-site agent network that runs on the existing
robotic-surgeries CI infrastructure. Cite claude-opus-47, claude-code,
claude-sonnet-46, ollama, and vllm.]

\subsection{Future Work to Reduce Approximations}
\label{sec:limits-future}

[LIMITATIONS SUBSECTION 5 PROMPT - PROCESS DURING THE PAPER POPULATION PASS]

[Document the concrete future work that reduces data approximations,
reduces surgery time below 1 minute, increases time resolution
beyond 1 kHz commands plus 10 kHz force, and integrates with the FDA
RTCT real-time signal-sharing framework. Enumerate at least the
following deliverables: (a) extend the 16-iteration sweep to 32 or
64 iterations at outputs/iterations/ to tighten the cumulative-arm-force
violation accounting; (b) add a competing 4-arm robot vendor entry
to the LLM tournament at outputs/comparison/ so the leaderboard is
not single vendor; (c) run an ablation study on the
cumulative-force-share clamp (the soft 11.0 N threshold) versus the
hard 12 N park threshold at outputs/iterations/; (d) increase the
sensor sample rate beyond 10 kHz force toward 100 kHz force per arm
to capture the high-frequency tip dynamics during Phase 2 bulk
resection; (e) compress the 60 second target to 30 seconds and to
15 seconds in follow-on variants that share this repository under
2030-gbm-30s/ and 2030-gbm-15s/; (f) ingest the FDA RTCT real-time
signal-sharing channel (cite fda2026realtime) and emit per-iteration
endpoints at hourly cadence; (g) adopt TRIPOD+AI and CREMLS reporting
standards (cite Collins2024TRIPODAI, ElEmam2024CREMLS) as the
reporting floor for any real-patient extension; (h) deposit the
release-aggregate L0 raw archive (currently 416 MB across 16
iterations) on Zenodo (cite zenodo) at DOI 10.5281/zenodo.18445179
with per-iteration pointer JSON resolution; (i) replace the
synthetic patient PAT-GBM-0001 with a TRIPOD+AI-compliant
retrospective cohort that pairs the 16-iteration sweep against real
EHR-linked oncology outcomes; (j) prototype both Track A and Track B
on the on-premises stack and report the per-iteration cost and
latency delta. State explicitly that this future work plan reduces
the approximations enumerated in \S\ref{sec:limits-accounting}, the
60-minute vs 1-minute deltas in \S\ref{sec:limits-deltas}, the
cross-simulation limitations in \S\ref{sec:limits-cross}, and the
Track A vs Track B trade offs in \S\ref{sec:limits-tracks}.]
